% !TEX root = ../nc-cubic.tex
\section{可換 Kuznetsov 実現を持つ幾何点の開性}
\label{grl:sec:geometric-locus}

主予想\ref{conj:nc-main}の「任意の対象」へ進むには、既知の一般点を増やすだけでなく、
可換模型の存在が族の中でどこまで保たれるかを調べる必要がある。
本節では、定理\ref{thm:nf-residual-formal}の通常の Artin 形式認識を、
有限表示の比較核のモジュライを通して代数的な開性へ移す。
四次元を除くという仮定と、実際の smooth proper な圏族という仮定を保つ。
この議論から境界での閉性や全内在的成分の被覆を結論しない。

\subsection{圏族と比較核の有限表示}

以下 $S$ は有限型 $\C$-スキームとする。
$S$ 上の圏族 $\mathcal C_S$ は、小さい冪等完備な $S$-線型 dg 圏の族で、
\'etale 局所的に
\[
 \mathcal C_S\simeq\operatorname{Perf}(\Gamma_S),\qquad
 \Gamma_S\in\operatorname{Perf}(S),\qquad
 \Gamma_S\in\operatorname{Perf}(\Gamma_S^{\mathrm{op}}
                      \otimes_S^{\mathbf L}\Gamma_S)
\]
と表示されるものとする。従って族は相対的に smooth proper である。
係数変更はこの表示の導来 tensor で定め、compact 対象とその Ind 完備化を
相対 Morita 理論で扱う。以下で核というときは perfect 双加群を意味する。
一般の Abel 圏族に、このような相対 dg 模型が自動的に存在するとは仮定しない。

$M_n\subset\PP(\operatorname{Sym}^3(\C^{n+2})^*)$ を
滑らかな埋込み三次超曲面の係数スキームとし、$\mathcal X\to M_n$ を普遍族とする。
相対例外列の射影核から $\operatorname{Ku}(\mathcal X/M_n)$ を構成する。
この圏族は smooth proper であり、\'etale 局所的な有限 dg 生成対象を持つ。

\begin{lemma}[可換比較を表示する核スタック]
\label{grl:lem:comparison-stack}
$S\times M_n$ 上に、対 $(X_T,P_T)$、すなわち
$M_n$ が与える三次族 $X_T$ と実際の Morita 同値
\[
 P_T:\mathcal C_T\xrightarrow{\sim}\operatorname{Ku}(X_T/T)
\]
をパラメータ付ける局所有限表示の Artin スタック $\mathcal E$ が存在する。
射は同じ埋込み三次族上の核の同型であり、核の scalar 自己同型は保持する。
幾何点の核の自己同型群は $\mathbf G_m$ である。
\end{lemma}

\begin{proof}
アファインな係数 chart 上で両側の smooth proper な有限 dg 模型を選ぶ。
二つの模型の導来 tensor 積も smooth proper であり、
To\"en--Vaqui\'e の対象モジュライ
\cite[Corollaries~3.15, 3.21, 3.22]{NFMToenVaquieModuli}
により、その perfect 双加群のうち負の自己 Ext がない部分は
局所有限表示の Artin スタックをなす。

perfect 双加群の左右随伴は相対双対で与えられる。
同値性は単位・余単位の二つの cone が零となる条件である。
これらの cone は底上 perfect なので、その零化 locus は開である
\cite[Lemma~3.25 and its proof]{NFMToenVaquieModuli}。
同補題は自己関手について記されているが、二つの dg 代数を用いる場合も、
各側の自由生成対象で単位・余単位を検査する同じ証明が適用される。
従って同値核の部分は開部分スタックである。

幾何点で同値核 $P$ を固定すれば、その自己 Hom は標的圏の
Hochschild cohomology である。命題\ref{prop:nf-hh}より
\[
 \operatorname{Ext}^{q}(P,P)=0\quad(q<0),\qquad
 \operatorname{End}(P)=\kappa .
\]
従って同値核は上の Artin 部分に属し、自己同型は $\kappa^\times$ となる。
核・随伴・単位・余単位と有限表示は相対 Morita 変更に自然であるため、
係数 chart 間で降下して $\mathcal E$ を与える。
ここでは核の同型類を先に選んで gerbe を中立化する操作は行わない。
\end{proof}

\subsection{形式認識から代数的な近傍への移行}

\begin{theorem}[四次元を除く可換実現 locus の開性]
\label{grl:thm:open-locus}
$n\ge3$、$n\ne4$ とし、上の smooth proper な族 $\mathcal C_S$ を取る。
幾何的なファイバーが滑らかな $n$ 次元三次超曲面の Kuznetsov 成分に
dg Morita 同値となる点の集合は、Zariski 開部分 $S_{\mathrm{cub}}\subset S$ である。
この開上では \emph{\'etale 局所的に}滑らかな三次族と実際の perfect 比較核が存在する。

より具体的に、$s\in S(\C)$ における
\[
 P_s:\mathcal C_s\xrightarrow{\sim}\operatorname{Ku}(X_s)
\]
を固定すると、$s$ の \'etale 近傍 $S'\to S$ 上で
$X_s$ と $P_s$ を、滑らかな三次族と Morita 同値へ延長できる。
この同値は全ての導来底変更に保たれる。
rotation、指定三次周期、構造対象および point scheme の標識は結論に含めない。
\end{theorem}

\begin{proof}
補題\ref{grl:lem:comparison-stack}の射 $q:\mathcal E\to S$ を考える。
まず $e=(X_s,P_s)$ が $\C$-点である場合に $q$ の形式的滑らかさを示す。
剰余体 $\C$ を持つ通常の局所 Artin 環の小拡大 $A'\twoheadrightarrow A$ と、
$\operatorname{Spec}A'\to S$、および $e$ を特殊ファイバーとする
$(X_A,P_A)\in\mathcal E(A)$ が与えられたとする。

特殊ファイバーの $P_s$ によって
$\operatorname{Ind}\mathcal C_{A'}$ を
$\operatorname{Ind}\operatorname{Ku}(X_s)$ の標識付き変形とみなす。
定理\ref{thm:nf-residual-formal}の本質的全射性から、
滑らかな三次の変形 $X_{A'}$ と、実際の圏同値の同型類
\[
 Q_{A'}:\mathcal C_{A'}\xrightarrow{\sim}\operatorname{Ku}(X_{A'})
\]
が得られる。compactly generated な圏の同値は compact 対象を保ち、
相対 Morita 理論によって perfect 双加群で表される。

ここでは存在だけでなく、既に指定された $P_A$ との整合が必要である。
$A$ 上では $P_AQ_{A'}^{-1}|_A$ が二つの幾何的変形の残余圏間の射を与える。
定理\ref{thm:nf-residual-formal}の\emph{群oidとしての完全忠実性}により、
この射は、特殊ファイバー上恒等な幾何的同型
$X_{A'}\times_{A'}A\simeq X_A$ によって実現される。
同定後には $Q_{A'}|_A\simeq P_A$ という核の実際の同型を選べる。
$\tau_{\le1}$ における射の等しさは、まさにこの自然同型の存在を意味する。
核の scalar 自己同型を自明と主張する必要はなく、
持上げに選んだ同型の非一意性も許す。

補題\ref{lem:nf-geometric-smooth}の最後の議論により、
$\OO_{X_A}(1)$ とその切断の基底を $A'$ へ持ち上げられる。
従って $X_{A'}$ を、指定された $M_n$ の $A$-点を持ち上げる
埋込み三次として選べる。これで $(X_A,P_A)$ の持上げを得た。

有限型の smooth atlas $U\to\mathcal E$ と $e$ の上の $\C$-点 $u$ を取る。
上の持上げと atlas の形式的滑らかさを合成すると、$U\to S$ は
$u$ において Artin 小拡大の持上げ条件を満たす。
局所有限表示と局所的な滑らかさの判定
\cite[Tag~02HX]{NFMStacksSmoothLocal}
より、$U\to S$ は $u$ で滑らかである。
この議論は $\mathcal E$ の全ての $\C$-点に適用できる。
有限型 $\C$-スキームは Jacobson であり、有限型 atlas の非滑らか locus が
非空なら $\C$-点を含むので、$q$ は全体で滑らかとなる。

従って $q$ の像は開であり、定義からその像が $S_{\mathrm{cub}}$ である。
smooth atlas を通し、滑らかな射の \'etale 局所切断
\cite[Tag~055U]{NFMStacksSmoothLocal}
を取れば、この開の \'etale 被覆上に実際の三次族と核が得られる。
固定した $\C$-点 $u$ を通る切断も取れる。
実際、$u$ の近傍の標準的な滑らか表示を
$U\to\mathbb A^d_S$ という \'etale 射にし、
$u$ の相対座標を持つ $S$ の切断に沿って引き戻せばよい。
これにより $X_s,P_s$ を保つ延長を得る。
最後に、核と逆核の単位・余単位の同型は導来 tensor に保たれるので、
構成した同値は全ての導来底変更後も同値である。
\end{proof}

\begin{remark}[主予想への寄与と残る障害]
\label{grl:rem:scope}
この定理は、既知の幾何的実現点を含む smooth proper な圏族の近傍についての結果である。
連結な $S$ の非空開で可換実現があることから、$S$ の全点での存在は従わない。
主予想へ進むには、境界の DVR に沿って三次模型と比較核を延長する条件、
および全内在的対象をこのような族で扱えることを別に示す必要がある。
系\paperref{moduli}{cor:nfm-residual-good-reduction}の選択した形式比較に依存する
良い還元も、比較核の境界延長を自動的には与えない。

四次元では命題\ref{prop:nf-fourfold-obstruction}の二つの余分な
Hochschild 変形方向があるため、この証明を適用できない。
rotation が実際の自己同値として持ち上がる場合は、
定理\ref{ifk:thm:rotation-formal}が固定形式空間を幾何的変形と同定する。
その結果、定理\ref{ifk:thm:rotation-neighborhood}では
相対 rotation と原点の rotation 付き核を入力として、四次元の局所比較を得る。
これは本定理の $n\ne4$ を無条件に除くものではない。
また本節では相対 smooth proper 性を入力とし、
任意の Abel 圏変形の飽和性や全 $D^b(\mathrm{coh})$ との一致を証明したとはしない。

形式認識と核モジュライの既知の代数性を組み合わせた帰結であり、
この開性の先行研究上の新規性の網羅的確認は未了である。
\end{remark}

\subsection{trait 上の比較核に残る二つの有限検査}

可換模型の滑らかな延長と、比較核の延長は別の条件である。
さらに、比較核を perfect 双加群として延長できても、
その特殊ファイバーが同値を与えるとは限らない。
以下は実際の延長核を得た後に残る検査を、有限複体として記述する。
核の存在をこの検査自体から導くものではない。

\begin{proposition}[DVR 上の同値核の有限判定]
\label{grl:prop:trait-kernel-test}
$R$ を DVR、$t$ をその一様化元、$K$ を商体、$k$ を剰余体とする。
$A,B$ を $R$ 上 smooth proper な dg 代数とし、
$P$ を $A$--$B$ perfect 双加群とする。
右加群の圏を用い、
\[
 F_P=-\otimes_A^{\mathbf L}P:
 \operatorname{Perf}(A)\longrightarrow\operatorname{Perf}(B)
\]
と置く。$F_{P_K}$ は同値であると仮定する。
$Q=\operatorname{RHom}_B(P,B)$ とし、単位・余単位の cone を
\[
 C_A=\operatorname{Cone}\left(
 A\longrightarrow\operatorname{RHom}_B(P,P)\right),\qquad
 C_B=\operatorname{Cone}\left(
 Q\otimes_A^{\mathbf L}P\longrightarrow B\right)
\]
とする。このとき以下が成立する。
\begin{enumerate}
\item $C_A,C_B$ は $R$ 上 perfect であり、その各 cohomology は有限長である。
\item 次の三条件は同値である：
  $F_P$ が同値であること、$F_{P_k}$ が同値であること、
  $(C_A\oplus C_B)\otimes_R^{\mathbf L}k=0$ であること。
\item 固定した $A,B,P$ に対して
\[
 d_R(P;A,B)=\sum_j\left(
 \operatorname{length}_R H^j(C_A)
 +\operatorname{length}_R H^j(C_B)\right)
\]
は有限な非負整数であり、上の条件は $d_R(P;A,B)=0$ と同値である。
\item $R\to R'$ を DVR の有限平坦な局所拡大、$e$ を分岐指数とすると、
\[
 d_{R'}(P_{R'};A_{R'},B_{R'})=e\,d_R(P;A,B).
\]
従って、固定した延長核の不良な特殊化は、有限分岐底変更だけでは解消しない。
\end{enumerate}
\end{proposition}

\begin{proof}
$P$ の $B$-perfect 性から、$F_P$ の右随伴は
$G_P=\operatorname{RHom}_B(P,-)\simeq-\otimes_B^{\mathbf L}Q$ で与えられる。
smooth proper 性と有限 cell--retract 表示により、
$P,Q$、二つの合成核、二つの cone は底上 perfect である。
随伴、単位・余単位、これらの cone の形成は導来底変更と両立する。

$A$ は $\operatorname{Perf}(A)$ の生成対象である。
従って $F_P$ の完全忠実性は $A$ 上で単位を検査すること、
すなわち $C_A=0$ と同値である。
同様に $B$ 上で余単位を検査することで、$C_A=C_B=0$ は
同値性と同値となる。
これは To\"en--Vaqui\'e
\cite[Lemma~3.25 and its proof]{NFMToenVaquieModuli}
の二つの代数に対する同じ生成対象の議論である。

一般ファイバーで同値であるため $(C_A)_K=(C_B)_K=0$ である。
従って有限生成な各 cohomology は $R$-torsion、従って有限長である。
底上 perfect な複体 $C$ について、導来 Nakayama により
$C\otimes_R^{\mathbf L}k=0$ なら $C=0$ である。
これを二つの cone に適用して第2項を得る。
第3項は有限長の非負性による。

最後に $R'$ の平坦性と底変更から
$H^j(C\otimes_R R')=H^j(C)\otimes_R R'$ である。
有限長 $R$-加群の組成列を用いると、
$\operatorname{length}_{R'}(M\otimes_R R')
=e\operatorname{length}_R M$ となる。
これを各 cohomology へ適用すれば第4項を得る。
\end{proof}

\begin{remark}[行列での検査とその依存性]
\label{grl:rem:trait-smith}
上の二つの cone を有界有限自由 $R$-複体で表せば、
特殊ファイバーでの判定は有限個の行列の $t=0$ での exactness の検査である。
DVR 上の Smith 標準形を逐次使うと、一般ファイバーで acyclic な
有限自由複体は、shift を除いて
$[R\xrightarrow{t^{m_\nu}}R]$、$m_\nu\ge0$、の有限直和に分解できる。
$m_\nu=0$ の項は contractible であり、残る指数の和が
上で定めた長さとなる。

この数値は固定した dg 生成模型と延長核の数値であり、
圏だけの新しい内在的不変量とはしない。
異なる延長核を選んだときの値の不変性も主張しない。
判定が有限であることと、具体的な族でその行列を既に計算したことは別である。
\end{remark}

\begin{example}[perfect 延長だけでは同値性は残らない]
\label{grl:ex:perfect-kernel-not-equivalence}
$A=B=R$ とし、$P=R\oplus R/(t)$ と置く。
$R/(t)$ は二項有限自由分解を持つので、$P$ は perfect 双加群である。
一般ファイバーでは $P_K=K$ で同値を与えるが、
特殊ファイバーでは
\[
 P\otimes_R^{\mathbf L}k\simeq k^{\oplus2}\oplus k[1]
\]
であり、tensor 関手は完全忠実でない。
この例では別の延長核 $R$ を選べば同値になる。
従って問題は有限底変更の有無だけでなく、適切な核の整格子の構成にある。
\end{example}

\begin{remark}[指定混合族の比較と選択済み一般点核の特殊化]
\label{grl:rem:mixed-trait-inputs}
指定された中心係数 $c$ の族については、
定理\paperref{mixed}{cmo:thm:all-central-commutative-kernel}が
滑らかな範囲全体で可換三次への実比較核を与える。
$c=0$ での直接構成は
定理\paperref{mixed}{cmo:thm:fixed-zero-commutative-kernel}である。
従って、各滑らかな $c$ における Ku 同値の存在を、
一般点核の特殊化に関する未解決問題として扱わない。

別に、非 PI 変形の $h$-族について選択済みの一般点核を固定した場合は、
定理\paperref{mixed}{mix:thm:h-exact-candidate}が一般点での実同値を、
定理\paperref{mixed}{mix:thm:h-integral-kernel}が perfect な整延長を与える。
この延長について残るのは、一般点での指定を保ちながら、
$h=1$ における unit と counit の cone がともに零となるように
整核を選べるかという点である。
これは特殊ファイバーにおける同値そのものの存在とは異なる条件であり、
smooth proper 性、判別式の一致、Hochschild 次元の一致だけからは従わない。
\end{remark}
